\chapter{局部凸空间}
\label{chap:locally-convex-spaces}

范数用一个非负数同时控制向量的所有方向。许多函数空间却有一系列自然的
控制量。例如，对 \(C^\infty[a,b]\) 中的函数，可能要求函数本身一致
接近，也可能同时要求前 \(m\) 阶导数一致接近；随着 \(m\) 改变，得到的
控制条件不能总由一个范数完整表达。局部凸空间以一族半范数代替单一范数，
同时保留线性运算、凸集和连续线性泛函所需的结构。

本章固定标量域 \(\mathbb K\in\{\R,\C\}\)。前半章从拓扑向量空间的
零点邻域出发，建立平衡集、吸收集、Minkowski 泛函和半范数族。后半章
讨论可度量性、完备性、弱拓扑与弱星拓扑。一般局部凸空间中用连续线性
泛函分离点或凸集所需的 Hahn--Banach 定理留到第四卷；本章只证明不依赖
该定理的有限维结论和弱拓扑的代数刻画。

\section{拓扑向量空间}

\begin{definition}[拓扑向量空间]
向量空间 \(E\) 若带有 Hausdorff 拓扑，并且加法
\[
 E\times E\to E,\qquad(x,y)\mapsto x+y
\]
与数乘
\[
 \mathbb K\times E\to E,\qquad(\lambda,x)\mapsto\lambda x
\]
连续，则称 \(E\) 为\term{拓扑向量空间}。
\end{definition}

每个平移 \(\tau_a(x)=a+x\) 都是同胚，逆映射为 \(\tau_{-a}\)。因此
拓扑完全由零向量的邻域决定，点 \(a\) 的邻域恰为 \(a+U\)，其中 \(U\)
遍历零点邻域。连续性还给出两条经常使用的缩小规则：对每个零点邻域
\(U\)，存在零点邻域 \(V\) 使
\[
 V+V\subseteq U,
\]
并可令 \(V=-V\)。

\begin{proposition}[线性映射只需在零点检验]
\label{prop:linear-map-continuity-zero}
设 \(E,F\) 为拓扑向量空间，\(T:E\to F\) 线性。下列条件等价：
\begin{enumerate}[label=\textup{(\roman*)}]
 \item \(T\) 连续；
 \item \(T\) 在 \(0\) 处连续；
 \item 对 \(F\) 的每个零点邻域 \(V\)，\(T^{-1}(V)\) 是 \(E\) 的
 零点邻域。
\end{enumerate}
\end{proposition}

\begin{proof}
前两项及第三项的等价是点连续性的原像表述。若 \(T\) 在零点连续，给定
\(a\in E\) 与 \(T(a)\) 的邻域 \(T(a)+V\)，可取零点邻域 \(U\) 使
\(T(U)\subseteq V\)。于是
\[
 T(a+U)=T(a)+T(U)\subseteq T(a)+V,
\]
故 \(T\) 在 \(a\) 处连续。
\end{proof}

拓扑向量空间作为加法群，具有
\chapterref{chap:uniform-spaces-topological-groups}所述的平移不变一致结构；
因加法交换，
左右一致结构相同。其基本周围可写成
\[
 \widetilde U=\{(x,y):x-y\in U\},
\]
其中 \(U\) 是零点邻域。于是 Cauchy 网的条件成为：对每个零点邻域
\(U\)，最终有 \(x_i-x_j\in U\)。

\section{平衡集、吸收集与凸邻域}

\begin{definition}[平衡、吸收与绝对凸]
设 \(A\subseteq E\)。
\begin{enumerate}[label=\textup{(\roman*)}]
 \item 若 \(|\lambda|\le1\) 蕴含 \(\lambda A\subseteq A\)，称 \(A\)
 \term{平衡}；
 \item 若对每个 \(x\in E\)，存在 \(t>0\) 使 \(x\in tA\)，称 \(A\)
 \term{吸收}；
 \item 若 \(A\) 同时凸且平衡，称 \(A\) \term{绝对凸}。
\end{enumerate}
\end{definition}

每个零点邻域都是吸收集：对固定 \(x\)，映射 \(\lambda\mapsto\lambda x\)
在 \(0\) 连续，故充分小的 \(\lambda\ne0\) 满足 \(\lambda x\in U\)，
即 \(x\in\lambda^{-1}U\)。平衡性则使从较小倍数回到集合内部的操作
合法。

\begin{proposition}[平衡邻域]
\label{prop:balanced-neighborhood-base}
拓扑向量空间中每个零点邻域都含有平衡的零点邻域。若空间还有由凸集组成
的零点邻域基，则每个零点邻域都含有绝对凸的零点邻域。
\end{proposition}

\begin{proof}
给定零点邻域 \(U\)，先在其中取一个开零点邻域并仍记为 \(U\)。数乘在紧集
\(\{\lambda:|\lambda|\le1\}\times\{0\}\) 的每一点连续。对每个
\(|\lambda|\le1\)，可取标量邻域 \(D_\lambda\) 与零点邻域
\(W_\lambda\)，使 \(D_\lambda W_\lambda\subseteq U\)。闭单位圆盘的
紧致性给出有限多个 \(D_{\lambda_k}\) 覆盖它。令
\(W=\bigcap_kW_{\lambda_k}\)，则 \(\mu W\subseteq U\) 对所有
\(|\mu|\le1\) 成立。集合
\[
 V=\bigcup_{|\mu|\le1}\mu W
\]
平衡、含 \(W\)，且包含于 \(U\)，所以是所需邻域。

若有凸邻域基，先在 \(U\) 内取凸邻域 \(U_0\)，再按上段取平衡邻域
\(V\subseteq U_0\)。\(V\) 的绝对凸包仍包含一个零点邻域，并因
\(U_0\) 凸而包含于 \(U_0\subseteq U\)。
\end{proof}

\begin{definition}[Minkowski 泛函]
若 \(A\subseteq E\) 吸收，定义
\[
 p_A(x)=\inf\{t>0:x\in tA\},
\]
称为 \(A\) 的\term{Minkowski 泛函}。
\end{definition}

\begin{theorem}[绝对凸集的 Minkowski 泛函]
\label{thm:minkowski-functional-seminorm}
若 \(A\) 吸收且绝对凸，则 \(p_A\) 是半范数，即
\[
 p_A(\lambda x)=|\lambda|p_A(x),\qquad
 p_A(x+y)\le p_A(x)+p_A(y).
\]
并且
\[
 \{x:p_A(x)<1\}\subseteq A\subseteq\{x:p_A(x)\le1\}.
\]
若 \(A\) 还是开集，则第一处包含为等号：
\(A=\{x:p_A(x)<1\}\)。
\end{theorem}

\begin{proof}
吸收性保证定义中的集合非空。对 \(\lambda\ne0\)，
\[
 \lambda x\in tA
 \quad\Longleftrightarrow\quad
 x\in(t/|\lambda|)A,
\]
其中使用平衡性处理复数相位，故齐次性成立；\(\lambda=0\) 直接。

若 \(x\in rA\)、\(y\in sA\)，其中 \(r,s>0\)，则
\[
 \frac{x+y}{r+s}
 =\frac r{r+s}\frac xr+\frac s{r+s}\frac ys\in A
\]
由凸性成立，因而 \(p_A(x+y)\le r+s\)。令
\(r\downarrow p_A(x)\)、\(s\downarrow p_A(y)\) 得三角不等式。

若 \(p_A(x)<1\)，取 \(p_A(x)<t<1\)。由下确界定义，存在 \(s<t\)
使 \(x\in sA\)，而平衡性给出 \(sA\subseteq A\)。若 \(x\in A\)，
则 \(p_A(x)\le1\)。最后设
\(A\) 开且 \(x\in A\)。映射 \(s\mapsto sx\) 在 \(s=1\) 连续，故
可取 \(s>1\) 使 \(sx\in A\)。于是 \(x\in s^{-1}A\)，
\(p_A(x)\le s^{-1}<1\)。
\end{proof}

\begin{theorem}[有限维拓扑向量空间]
\label{thm:finite-dimensional-tvs-euclidean}
每个有限维 Hausdorff 拓扑向量空间都线性同胚于
\(\mathbb K^n\) 的通常拓扑。特别地，它的所有线性泛函都连续。
\end{theorem}

\begin{proof}
取基 \(e_1,\ldots,e_n\)，线性双射
\[
 T:\mathbb K^n\to E,\qquad T(a_1,\ldots,a_n)=\sum_{k=1}^na_ke_k
\]
由有限次加法与数乘的连续性而连续。只需证 \(T^{-1}\) 在零点连续。

令 \(S\) 为 \(\mathbb K^n\) 的欧氏单位球面。\(T(S)\) 紧，而 \(E\)
Hausdorff，所以 \(T(S)\) 闭且不含零点。取平衡零点邻域
\(V\subseteq E\setminus T(S)\)。若 \(T(a)\in V\) 且
\(\norm{a}_2\ge1\)，则 \(a=rs\)，其中 \(r\ge1\)、\(s\in S\)。平衡性
给出 \(T(s)=r^{-1}T(a)\in V\)，与 \(V\cap T(S)=\varnothing\) 矛盾。
因此 \(T^{-1}(V)\) 包含于欧氏单位球；对任意 \(\varepsilon>0\)，
\(T^{-1}(\varepsilon V)\) 包含于半径 \(\varepsilon\) 的球。这证明逆映射
在零点连续。

任意线性泛函与 \(T\) 复合后是 \(\mathbb K^n\) 上的坐标线性组合，
所以连续。
\end{proof}

\section{半范数族}

\begin{definition}[半范数]
映射 \(p:E\to[0,+\infty)\) 若满足
\[
 p(\lambda x)=|\lambda|p(x),\qquad
 p(x+y)\le p(x)+p(y),
\]
则称为\term{半范数}。半范数允许 \(p(x)=0\) 而 \(x\ne0\)。
\end{definition}

设 \(\mathcal P\) 是 \(E\) 上的一族半范数。有限多个半范数和
\(\varepsilon>0\) 给出集合
\[
 U(p_1,\ldots,p_n;\varepsilon)
 =\{x\in E:p_k(x)<\varepsilon, 1\le k\le n\}.
\]

\begin{theorem}[半范数族生成的拓扑]
\label{thm:seminorm-generated-topology}
上述集合构成一个向量拓扑的零点邻域基。该拓扑是使
\(\mathcal P\) 中每个半范数连续的最弱向量拓扑，并且它为 Hausdorff
拓扑当且仅当 \(\mathcal P\) 分离点，即
\[
 \bigcap_{p\in\mathcal P}\ker p=\{0\}.
\]
只有在这一分离条件成立时，它才按本书的 Hausdorff 约定成为拓扑向量
空间。
\end{theorem}

\begin{proof}
有限交仍可通过合并半范数得到更小的同类集合；三角不等式给出
\[
 U(p_1,\ldots,p_n;\varepsilon/2)
 +U(p_1,\ldots,p_n;\varepsilon/2)
 \subseteq U(p_1,\ldots,p_n;\varepsilon),
\]
齐次性则给出数乘在零点附近的连续性。平移这些零点邻域便得到向量拓扑。
每个 \(p\in\mathcal P\) 连续，因为 \(\{x:p(x)<\varepsilon\}\) 是零点
邻域；任何使全部 \(p\) 连续的向量拓扑都必须包含这些有限交，故当前
拓扑最弱。

若非零 \(x\) 被某个 \(p\) 检出，即 \(p(x)>0\)，则
\(p\)-小球可把 \(0,x\) 分开。若 \(x\ne0\) 属于全部核，则它属于每个
零点基本邻域，拓扑不可能 Hausdorff。
\end{proof}

\begin{proposition}[半范数空间中的连续线性映射]
\label{prop:linear-map-seminorm-estimate}
设 \(E,F\) 的拓扑分别由半范数族 \(\mathcal P,\mathcal Q\) 生成。
线性映射 \(T:E\to F\) 连续，当且仅当对每个 \(q\in\mathcal Q\)，
存在 \(p_1,\ldots,p_n\in\mathcal P\) 和 \(C>0\)，使
\[
 q(Tx)\le C\max_{1\le k\le n}p_k(x)\qquad(x\in E).
\]
\end{proposition}

\begin{proof}
若上述估计成立，则对任意 \(\varepsilon>0\)，只要
\(\max_k p_k(x)<\varepsilon/C\)，便有 \(q(Tx)<\varepsilon\)，故
\(T\) 在零点连续。反之，若 \(T\) 连续，则对集合
\(\{y:q(y)<1\}\) 的原像，可取 \(\varepsilon>0\) 和
\(p_1,\ldots,p_n\)，使
\[
 \max_kp_k(x)<\varepsilon\quad\Longrightarrow\quad q(Tx)<1.
\]
记 \(r=\max_kp_k(x)\)。若 \(r>0\)，对
\(\lambda=\varepsilon/(2r)\) 应用上式，得
\(q(Tx)<2r/\varepsilon\)。若 \(r=0\)，则对任意 \(t>0\) 都可把
\(tx\) 代入，齐次性迫使 \(q(Tx)=0\)。故可取 \(C=2/\varepsilon\)。
\end{proof}

\begin{example}[乘积空间]
在 \(\mathbb K^S\) 上，对每个 \(s\in S\) 令
\(p_s(x)=|x(s)|\)。坐标半范数族生成的拓扑正是乘积拓扑。它允许在一个
基本邻域中只限制有限多个坐标。
\end{example}

\begin{example}[光滑函数空间]
在 \(C^\infty[a,b]\) 上定义
\[
 p_m(f)=\max_{0\le k\le m}\norm{f^{(k)}}_\infty,
 \qquad m\in\N.
\]
这些半范数分离点。相应拓扑中的 \(f_n\to f\) 表示对每个固定
\(m\)，前 \(m\) 阶导数都一致收敛；这里的 \(m\) 可以随所考察的邻域
改变，不要求存在一个同时控制全部阶导数的有限常数。
\end{example}

\section{局部凸拓扑}

\begin{definition}[局部凸空间]
若拓扑向量空间 \(E\) 的零点具有由凸集组成的邻域基，则称 \(E\) 为
\term{局部凸空间}。
\end{definition}

半范数球是绝对凸集，所以半范数族生成的 Hausdorff 拓扑局部凸。反向也
成立。

\begin{theorem}[局部凸拓扑的半范数表示]
\label{thm:locally-convex-seminorm-representation}
Hausdorff 拓扑向量空间局部凸，当且仅当其拓扑由一族分离点的连续半范数
生成。
\end{theorem}

\begin{proof}
只需证明正向。由命题 \ref{prop:balanced-neighborhood-base}，每个零点邻域
都含有一个绝对凸的零点邻域；再取其内部，便得到由开绝对凸集组成的
零点邻域基 \(\mathcal B\)。对每个 \(U\in\mathcal B\)，
Minkowski 泛函 \(p_U\) 是半范数，并且
\[
 U=\{x:p_U(x)<1\}.
\]
所以 \(p_U\) 连续，且这些单位球恰好恢复原零点邻域基。Hausdorff 性说明
任意 \(x\ne0\) 都被某个零点邻域排除，因而某个 \(p_U(x)>0\)；半范数
族分离点。
\end{proof}

为了判断一个局部凸拓扑能否来自单一范数，需要一个不依赖范数的有界概念。

\begin{definition}[拓扑有界集]
集合 \(B\subseteq E\) 称为\term{有界}，若对每个零点邻域 \(U\)，
存在 \(t>0\) 使 \(B\subseteq tU\)。
\end{definition}

\begin{theorem}[局部凸空间的赋范化判据]
\label{thm:locally-convex-normability}
Hausdorff 局部凸空间的拓扑可由一个范数生成，当且仅当零点有一个有界
邻域。
\end{theorem}

\begin{proof}
范数拓扑的单位球是有界零点邻域。反之，设 \(B\) 是有界零点邻域。在
\(B\) 内取开绝对凸邻域 \(U\)，则 \(U\) 仍有界。其 Minkowski 泛函
\(p_U\) 是半范数。若 \(p_U(x)=0\) 且 \(x\ne0\)，取平衡零点邻域
\(V\) 使 \(x\notin V\)。有界性给出 \(U\subseteq rV\)。而
\(p_U(x)=0\) 蕴含 \(x\in r^{-1}U\subseteq V\)，矛盾。因此
\(p_U\) 是范数。

集合 \(U=\{p_U<1\}\) 是原拓扑中的开邻域，所以每个范数球都是原拓扑
开集。对任意原零点邻域 \(W\)，有界性给出 \(U\subseteq tW\)，即
\(t^{-1}U\subseteq W\)，所以每个原邻域又含有一个范数球。两种拓扑相同。
\end{proof}

若 \(S\) 无限，乘积空间 \(\mathbb K^S\) 没有有界零点邻域。事实上，
任一基本邻域只限制有限坐标；选一个未受限制的坐标 \(s\)，该邻域便包含
整条直线 \(\{te_s:t\in\mathbb K\}\)，不可能被固定倍数的
\(\{|x(s)|<1\}\) 吸收。因此乘积拓扑局部凸，却不可赋范化。

\section{可度量化}

\begin{theorem}[局部凸空间的可度量化]
\label{thm:locally-convex-metrizable}
Hausdorff 局部凸空间可度量化，当且仅当其拓扑可由可数族半范数
\((p_n)_{n\in\Npos}\) 生成。此时可取平移不变度量
\[
 d(x,y)=\sum_{n=1}^\infty2^{-n}
 \frac{p_n(x-y)}{1+p_n(x-y)}.
\]
\end{theorem}

\begin{proof}
若有可数生成族，级数一致收敛且每项非负。分离点保证 \(d(x,y)=0\)
只在 \(x=y\) 时成立。令 \(h(t)=t/(1+t)\)。对 \(s,t\ge0\)，有
\[
 h(s)+h(t)-h(s+t)
 =\frac{st(2+s+t)}{(1+s)(1+t)(1+s+t)}\ge0.
\]
函数 \(h\) 又单调递增，故半范数三角不等式给出
\[
 h\bigl(p_n(x-z)\bigr)
 \le h\bigl(p_n(x-y)\bigr)+h\bigl(p_n(y-z)\bigr).
\]
逐项乘以 \(2^{-n}\) 并求和，得到 \(d\) 的三角不等式。由于
\(0\le h\le1\)，定义 \(d\) 的级数在 \(E\times E\) 上一致收敛。

下面比较两个零点邻域系。给定 \(d\)-球 \(B_d(0,r)\)，先取 \(N\) 使
\[
 \sum_{n>N}2^{-n}<\frac r2.
\]
再对 \(1\le n\le N\) 选取 \(\eta_n>0\)，使
\[
 \sum_{n=1}^N2^{-n}h(\eta_n)<\frac r2.
\]
于是有限半范数邻域
\[
 U=\{x:p_n(x)<\eta_n,\ 1\le n\le N\}
\]
包含于 \(B_d(0,r)\)。反过来，给定基本邻域
\[
 W=\{x:p_{n_k}(x)<\varepsilon_k,\ 1\le k\le m\},
\]
取
\[
 0<\delta<\min_{1\le k\le m}
 2^{-n_k}h(\varepsilon_k).
\]
若 \(d(x,0)<\delta\)，则对每个 \(k\)，
\[
 2^{-n_k}h\bigl(p_{n_k}(x)\bigr)\le d(x,0)
 <2^{-n_k}h(\varepsilon_k),
\]
由 \(h\) 严格递增可知 \(p_{n_k}(x)<\varepsilon_k\)。所以
\(B_d(0,\delta)\subseteq W\)，两种拓扑相同。

反之，若空间可度量化，则零点第一可数。取递减的可数零点邻域基
\((U_n)\)，并在每个 \(U_n\) 内取开绝对凸邻域 \(V_n\)。Minkowski
半范数 \(p_{V_n}\) 的单位球就是 \(V_n\)，所以这个可数半范数族生成
原拓扑。
\end{proof}

可度量化不等于可赋范化。\(C^\infty[a,b]\) 由可数半范数 \(p_m\)
生成，因而可度量化；但任何零点邻域只限制有限阶导数，高阶方向仍可任意
放大，所以没有有界零点邻域，不能由一个范数生成。

\section{完备局部凸空间}

局部凸空间采用其加法群的一致结构。网 \((x_i)\) 为 Cauchy 网，当且仅当
对每个连续半范数 \(p\) 和 \(\varepsilon>0\)，最终有
\[
 p(x_i-x_j)<\varepsilon.
\]
在可度量化情形，只需检验 Cauchy 列；一般情形仍应使用网或滤子。

\begin{definition}[Fréchet 空间]
完备、可度量化的局部凸空间称为\term{Fréchet 空间}。
\end{definition}

\begin{proposition}[乘积的完备性]
\label{prop:product-locally-convex-complete}
任意一族完备局部凸空间的乘积在乘积拓扑下完备。特别地，
\(\mathbb K^S\) 对任意集合 \(S\) 都完备。
\end{proposition}

\begin{proof}
设 \(\mathcal F\) 是乘积中的 Cauchy 滤子。每个坐标投影一致连续，所以
像滤子在相应因子中 Cauchy，因而收敛到某个 \(x_s\)。令
\(x=(x_s)_{s\in S}\)。乘积基本邻域只限制有限多个坐标；把这些坐标上
的最终条件作有限交，便得到 \(\mathcal F\to x\)。
\end{proof}

\begin{theorem}[光滑函数空间的完备性]
\label{thm:smooth-function-frechet}
设 \(a<b\)。空间 \(C^\infty[a,b]\) 配备半范数
\[
 p_m(f)=\max_{0\le k\le m}\norm{f^{(k)}}_\infty
\]
时是 Fréchet 空间。
\end{theorem}

\begin{proof}
可数半范数族给出可度量化。设 \((f_n)\) 是 Cauchy 列。对每个
\(k\in\N\)，\((f_n^{(k)})\) 在 \(C[a,b]\) 的上确界范数下 Cauchy，
故一致收敛到某个连续函数 \(g_k\)。

\chapterref{chap:uniform-convergence-preservation}的逐项求导定理应用于
函数列 \((f_n^{(k)})\)：其在一点的值收敛，导函数
\(f_n^{(k+1)}\) 一致收敛到 \(g_{k+1}\)，所以
\(g_k\) 可导且 \(g_k'=g_{k+1}\)。从 \(k=0\) 递推，得到
\(g_0\in C^\infty[a,b]\) 且 \(g_0^{(k)}=g_k\)。每个固定 \(m\) 下
前 \(m\) 阶一致收敛，故 \(p_m(f_n-g_0)\to0\)。所以空间完备。
\end{proof}

完备性依赖所取拓扑。同一个向量空间换用更弱的半范数族，Cauchy 对象会
增多；换用更强的半范数族，收敛要求会增加。因而写“函数空间完备”时，
必须同时给出其拓扑。

\section{弱拓扑和弱星拓扑}

设 \(E\) 为向量空间，\(F\subseteq E^\#\) 为代数对偶中的线性子空间；
这里 \(E^\#\) 表示全部线性泛函。对 \(\varphi\in F\)，令
\(p_\varphi(x)=|\varphi(x)|\)。

\begin{definition}[弱拓扑]
由半范数族 \(\{p_\varphi:\varphi\in F\}\) 生成的拓扑称为
\(E\) 关于 \(F\) 的\term{弱拓扑}，记作 \(\sigma(E,F)\)。它是使
\(F\) 中全部线性泛函连续的最弱向量拓扑。
\end{definition}

零点的基本弱邻域为
\[
 U(\varphi_1,\ldots,\varphi_n;\varepsilon)
 =\{x:|\varphi_k(x)|<\varepsilon, 1\le k\le n\}.
\]
因而网 \(x_i\) 在 \(\sigma(E,F)\) 中趋于 \(x\)，当且仅当
\[
 \varphi(x_i)\to\varphi(x)\qquad(\varphi\in F).
\]
该拓扑 Hausdorff 当且仅当 \(F\) 分离 \(E\) 中的点。
若 \(F\) 不分离点，\(\sigma(E,F)\) 仍是有用的线性拓扑，但不符合本书
对拓扑向量空间所作的 Hausdorff 约定。

若 \(E\) 已是拓扑向量空间，记其连续对偶为
\[
 E'=\{\varphi:E\to\mathbb K:\varphi\text{ 线性且连续}\}.
\]
此时 \(\sigma(E,E')\) 简称 \(E\) 的弱拓扑。它不强于原拓扑；原拓扑
的收敛必推出弱收敛，逆向通常不成立。

\begin{definition}[弱星拓扑]
在 \(E'\) 上，由评价泛函
\[
 \operatorname{ev}_x:E'\to\mathbb K,\qquad
 \operatorname{ev}_x(\varphi)=\varphi(x),\qquad x\in E,
\]
生成的拓扑称为\term{弱星拓扑}，记作 \(\sigma(E',E)\)。
\end{definition}

网 \(\varphi_i\) 弱星收敛到 \(\varphi\)，当且仅当对每个固定
\(x\in E\)，\(\varphi_i(x)\to\varphi(x)\)。不同线性泛函必在某个
向量上取值不同，所以弱星拓扑总是 Hausdorff。第四卷的
Banach--Alaoglu 定理将把弱星拓扑解释为某个乘积空间子空间上的拓扑，
并用 Tychonoff 定理得到紧致性。

\section{连续线性泛函的分离作用}

弱拓扑的连续泛函没有隐藏的新成员；连续性迫使它只依赖有限多个定义弱
拓扑的坐标。

\begin{theorem}[弱拓扑的连续对偶]
\label{thm:dual-of-weak-topology}
设 \(F\subseteq E^\#\) 为线性子空间。在 \(E\) 配备
\(\sigma(E,F)\) 后，其连续对偶恰为 \(F\)。
\end{theorem}

\begin{proof}
按定义，\(F\) 中每个泛函连续。反过来，设 \(g:E\to\mathbb K\) 在弱
拓扑下连续。存在 \(\varphi_1,\ldots,\varphi_n\in F\) 和
\(\varepsilon>0\)，使
\[
 |\varphi_k(x)|<\varepsilon\ (1\le k\le n)
 \quad\Longrightarrow\quad |g(x)|<1.
\]
若 \(x\in\bigcap_k\ker\varphi_k\)，则对任意 \(t>0\)，\(tx\) 仍满足
左侧条件，于是 \(t|g(x)|<1\)。令 \(t\to\infty\)，得 \(g(x)=0\)。
所以
\[
 \bigcap_{k=1}^n\ker\varphi_k\subseteq\ker g.
\]
定义 \(T:E\to\mathbb K^n\) 为
\(T(x)=(\varphi_1(x),\ldots,\varphi_n(x))\)。上式保证
可以在子空间 \(T(E)\) 上定义线性泛函
\[
 \widetilde g(Tx)=g(x).
\]
事实上，若 \(Tx=Ty\)，则 \(x-y\in\bigcap_k\ker\varphi_k\subseteq
\ker g\)，所以右端与代表元无关。
有限维线性代数把 \(\widetilde g\) 延拓到 \(\mathbb K^n\)，故存在
\(a_1,\ldots,a_n\) 使
\[
 g(x)=\sum_{k=1}^na_k\varphi_k(x).
\]
由于 \(F\) 是线性子空间，\(g\in F\)。
\end{proof}

\begin{corollary}[点分离与弱拓扑]
\label{cor:weak-topology-separation}
\(\sigma(E,F)\) 为 Hausdorff 拓扑，当且仅当 \(F\) 分离点。若成立，
映射
\[
 J:E\to\mathbb K^F,\qquad J(x)=(\varphi(x))_{\varphi\in F}
\]
是线性单射，且弱拓扑正是由乘积拓扑在 \(J(E)\) 上诱导的拓扑。
\end{corollary}

\begin{proof}
Hausdorff 判据已由半范数族的核交给出。若 \(F\) 分离点，\(J\) 单射；
乘积拓扑由有限多个坐标控制，拉回后的基本邻域正是弱基本邻域。因此
\(J\) 是到其像的线性同胚。
\end{proof}

有限维 Hausdorff 拓扑向量空间的连续对偶由坐标泛函组成，当然分离点。
对一般 Hausdorff 局部凸空间，结论仍成立，但证明需要把某个非零向量上
定义的线性泛函延拓到全空间，同时保持一个 Minkowski 半范数控制；这正是
Hahn--Banach 定理承担的工作。后文将证明：若 \(x\ne0\)，存在
\(f\in E'\) 使 \(f(x)\ne0\)；更强地，点与适当闭凸集可由连续线性
泛函分离。在该定理建立以前，本章不把这一结论当作已知事实使用。

\section*{习题}
\addcontentsline{toc}{section}{习题}

\subsection*{A组　邻域、半范数与连续性}
\begin{enumerate}[label=\textbf{\thechapter.\arabic*.}]
 \exerciseitem{tvs-translations-homeomorphisms}
 证明拓扑向量空间中的平移和非零数乘都是同胚，并由此说明零点邻域确定
 全部拓扑。
 \exerciseitem{balanced-absorbing-neighborhood}
 证明每个零点邻域吸收；给出吸收但不平衡、平衡但不凸、凸但不平衡的
 \(\R^2\) 中例子。
 \exerciseitem{minkowski-open-unit-ball}
 补全定理 \ref{thm:minkowski-functional-seminorm} 中复数齐次性和
 \(A=\{p_A<1\}\) 的证明，并说明开性为何不能随意删除。
 \exerciseitem{seminorm-quotient-norm}
 设 \(p\) 为半范数。证明 \(E/\ker p\) 上
 \(\norm{[x]}=p(x)\) 良定义且为范数。
 \exerciseitem{linear-map-seminorm-bound}
 对命题 \ref{prop:linear-map-seminorm-estimate}，证明用
 \(p_1+\cdots+p_n\) 替代 \(\max p_k\) 得到等价判据。
 \exerciseitem{finite-dimensional-topology-uniqueness}
 证明有限维 Hausdorff 拓扑向量空间的任意线性子空间闭，任意线性映射到
 另一个有限维拓扑向量空间都连续。
\end{enumerate}

\subsection*{B组　局部凸、可度量与完备}
\begin{enumerate}[label=\textbf{\thechapter.\arabic*.},resume]
 \exerciseitem{product-seminorm-topology}
 证明坐标半范数在 \(\mathbb K^S\) 上生成乘积拓扑，并证明当 \(S\)
 无限时该空间不可赋范化。
 \exerciseitem{countable-seminorm-metric-details}
 验证
 \[
 d(x,y)=\sum_{n\ge1}2^{-n}\frac{p_n(x-y)}{1+p_n(x-y)}
 \]
 的三角不等式，并证明它与半范数族生成同一拓扑。
 \exerciseitem{bounded-set-seminorm-characterization}
 设拓扑由半范数族 \(\mathcal P\) 生成。证明 \(B\subseteq E\) 有界，
 当且仅当
 \[
  \sup_{x\in B}p(x)<\infty\qquad(p\in\mathcal P).
 \]
 \exerciseitem{smooth-space-not-normable}
 证明 \(C^\infty[a,b]\) 在半范数 \(p_m\) 下没有有界零点邻域，因而
 不可赋范化。
 \exerciseitem{product-completeness-coordinates}
 用 Cauchy 网而非滤子重新证明任意完备局部凸空间的乘积完备。
 \exerciseitem{smooth-frechet-completeness}
 逐项核对光滑函数空间完备性证明中
 \chapterref{chap:uniform-convergence-preservation}逐项求导定理的假设，
 并证明极限确在每个 \(p_m\) 下收敛。
\end{enumerate}

\subsection*{C组　弱拓扑与连续泛函}
\begin{enumerate}[label=\textbf{\thechapter.\arabic*.},resume]
 \exerciseitem{weak-convergence-net-criterion}
 从基本邻域定义证明：\(x_i\to x\) 于 \(\sigma(E,F)\) 当且仅当
 \(\varphi(x_i)\to\varphi(x)\) 对每个 \(\varphi\in F\) 成立。
 \exerciseitem{weak-dual-finite-coordinate-factorization}
 不引用定理 \ref{thm:dual-of-weak-topology} 的结论，完整重建其中的有限
 坐标分解，并说明 \(F\) 必须取线性子空间。
 \exerciseitem{weak-star-product-embedding}
 证明弱星拓扑是映射
 \[
 E'\to\mathbb K^E,\qquad\varphi\mapsto(\varphi(x))_{x\in E}
 \]
 拉回的乘积拓扑，并据此证明它 Hausdorff。
 \exerciseitem{weak-topology-finite-dimensional}
 设 \(E\) 有限维且 \(F\subseteq E^\#\) 分离点。证明
 \(\sigma(E,F)\) 等于 \(E\) 的通常有限维拓扑。
 \optionalexerciseitem{canonical-map-second-dual-injectivity}
 对拓扑向量空间 \(E\)，定义
 \(J:E\to(E')^\#\) 为 \(J(x)(f)=f(x)\)。证明 \(J\) 单射当且仅当
 \(E'\) 分离点；解释为何一般 Hausdorff 局部凸空间中的单射性需等到
 Hahn--Banach 定理后才能使用。
\end{enumerate}
